\section{Purpose of laboratory work}

The purpose of this laboratory work is to implement an analog signal acquisition and
conditioning system for a digital temperature sensor on the ESP32 microcontroller,
using FreeRTOS for concurrent task execution and an OLED display plus STDIO for reporting.

\subsection{Objectives of the laboratory work}

Laboratory work has the following objectives that should be achieved:

\begin{itemize}
  \item Implementing periodic sensor acquisition from a DHT11 digital temperature
    sensor at a configurable recurrence (2000\,ms), exposing
    the raw value through an internal \texttt{sensor\_read()} interface, with
    support for selecting a sensor from a catalogue of available sensors
  \item Implementing analog signal conditioning through a three-stage preprocessing
    pipeline: saturation (clamping to the sensor's physical range $0$--$50\,^\circ$C),
    a median filter (window size~5) for salt-and-pepper noise removal, and an
    exponential moving average (EMA, $\alpha = 0.2$) for weighted smoothing
  \item Evaluating alert thresholds on the processed signal with hysteresis
    (ON $\geq 28\,^\circ$C, OFF $\leq 24\,^\circ$C) to generate a stable binary
    alert state that drives LED indicators
  \item Displaying a structured report at a lower recurrence (500\,ms) on an OLED
    128$\times$64 display and via STDIO, showing the raw, median-filtered, and
    EMA-averaged temperature values, sensor validity, and alert state
  \item Structuring the firmware into three independent FreeRTOS tasks (Acquisition,
    Conditioning, Reporter) communicating through a binary semaphore and a mutex
  \item Organizing the codebase into three hardware-abstraction layers: MCAL, ECAL, and SRV
\end{itemize}

\subsection{Problem definition}

The application must be structured in three distinct FreeRTOS tasks:

\begin{enumerate}
  \item \textbf{Task 1 -- Sensor Acquisition (TaskAcquisition):} Reads raw temperature
    data from the currently selected sensor (DHT11) every 2000\,ms. The active sensor
    is chosen at startup from a sensor catalogue. The raw value is stored in shared
    sensor data and a binary semaphore is given to signal the conditioning task.

  \item \textbf{Task 2 -- Signal Conditioning (TaskConditioning):} Blocked on the binary
    semaphore; wakes immediately when new sensor data arrives. Applies a three-stage
    preprocessing pipeline:
    \begin{enumerate}
      \item \textbf{Saturation} -- clamps the raw value to the physical range of the
        sensor ($0$--$50\,^\circ$C for DHT11)
      \item \textbf{Median filter} (window = 5) -- removes salt-and-pepper outliers by
        sorting the last 5 saturated samples and selecting the median
      \item \textbf{Exponential Moving Average} ($\alpha = 0.2$) -- applies weighted
        smoothing: $y_n = \alpha \cdot x_n + (1 - \alpha) \cdot y_{n-1}$
    \end{enumerate}
    After filtering, the task evaluates alert thresholds with hysteresis
    (ON $\geq 28\,^\circ$C, OFF $\leq 24\,^\circ$C) and updates LED indicators:
    green LED when temperature is within normal range, red LED when the alert is active.

  \item \textbf{Task 3 -- Display and Reporting (TaskReporter):} Runs every 500\,ms
    using \texttt{vTaskDelayUntil} for drift-corrected periodicity. Updates the OLED
    128$\times$64 display with all intermediate pipeline values (raw, median-filtered,
    EMA-averaged), sensor name, and alert state. Emits a serial plotter line in the
    \texttt{>name:value} format and every 5\,seconds prints a full structured report
    including pipeline parameters and threshold configuration.
\end{enumerate}
